% ===================================================================
%  TABELL Nr. 7 — D'Duerf am Wanter. De Bauerenhaff am Fréijoer.
%  Traduction 2026 d'apres texto/io, controlee sur texto/fr, texto/de
%  et texto/nl. Consigne : voir texto/lb/10-tabell-01.tex.
%
%  LE VILLAGE DE LA PREMIERE SCENE ALIGNE SIX METIERS, ET UN SEUL
%  PORTE UN NOM FRANCAIS. Le forgeron est « Schmadd », le cordonnier
%  « Schouster », le charron « Stellmécher », le bourrelier
%  « Suedler », le meunier « Millermeeschter » — tous germaniques. Le
%  barbier, lui, est « Coiffeur ».
%
%  ET CE N'EST PAS LA LIGNE DU TABLEAU 5, C'EST UNE LIGNE PLUS FINE.
%  Le tableau 5 opposait le travail et la societe, l'atelier et le
%  salon. Ici les six hommes travaillent tous de leurs mains, dans la
%  meme rue, pour les memes clients. Ce qui separe le sixieme des
%  cinq autres, c'est que LES CINQ FONT OU REPARENT UNE CHOSE ET QUE
%  LE SIXIEME S'OCCUPE D'UNE PERSONNE. Le fer, la chaussure, la roue,
%  le harnais, la farine ont des noms allemands ; la coiffure a un
%  nom francais.
%
%  LA BLAGUE DU CORDONNIER TIENT, comme elle tenait en lituanien. Le
%  fac-simile oppose « shuifisto » et « shureparisto » ; le
%  luxembourgeois a « Schouster » celui qui fait les chaussures et
%  « Schouflécker » celui qui les rapetasse, batis sur le meme mot,
%  de sorte que l'opposition se lit sans glose.
%
%  ET LA FERME NE COUTE PAS UN EMPRUNT, comme en tcheque et en
%  lituanien. « Kou » la vache, « Stéier » le taureau, « Kallef » le
%  veau, « Schof » la brebis, « Hämmel » le belier, « Lamm »
%  l'agneau, « Geess » la chevre, « Sau » la truie, « Ferkel » le
%  porcelet, « Mär » la jument, « Fillen » le poulain, « Hong » la
%  poule, « Hunn » le coq, « Int » le canard, « Gäns » l'oie,
%  « Dauf » le pigeon, « Kanéngchen » le lapin. Trois colonnes, trois
%  familles de langues, et le meme resultat au meme tableau : LA
%  FERME EST L'ENDROIT OU AUCUNE LANGUE N'A EU BESOIN D'EMPRUNTER.
% ===================================================================

%%K t07-noto-1 noto
\VUnotes{143.2mm}{%
(*) Fir d'Eenzelheeten iwwert d'Iessen, kuck d'Tabellen 6 an 13.%
}

%%K t07-apar-1 apar
\VUsaut{4.0mm}
\VUcentre{13.2pt}{90}{ZWEET SERIE}
\VUsaut{3.0mm}
\VUfilet{16mm}
\VUsaut{5.6mm}
\VUtitre{15.0pt}{60}{TABELL Nr. 7}
\VUsaut{3.0mm}

%%K t07-tit-1 sub
\VUcentre{12.0pt}{0}{\textit{Éischt Zeen.}}
\VUsaut{3.0mm}

%%K t07-tit-2 sub
\VUpk{10.2pt}{40}{D'Duerf am Wanter.}
\VUsaut{4.0mm}

%%K t07-c1-01-1 p
1. --- An de Neijoersvakanzen hunn ech mäi Papp op Neistad begleet, an e Stiedchen,
wou hien e puer Handelsaffären ze reegele hat.

%%K t07-c1-02-1 p
2. --- Mir hunn d'\VUgras{Diligence} \textsuperscript{(11)} geholl, déi tëscht der
Stad an deem Fleck verkéiert; mä well et ganz kal war, war déi Rees an engem wéineg
bequeme Gefier net ganz agreabel.

%%K t07-c1-03-1 p
3. --- Dofir hu mir e richtege Genoss gespuert, wéi mir de \VUgras{Kutscher} säi Won
um Plaz virun der \VUgras{Auberge} \textsuperscript{(2)} vum \textit{Wäisse Bier}
ustoppe gesinn hunn. Den \VUgras{Aubergiste} \textsuperscript{(4)} huet op eis op der
Schwell vu sengem Haus gewaart a rouheg seng \VUgras{Päif} gefëmmt.

%%K t07-c1-03-2 p
Hien huet eis erageruff, an ech hu mech net laang biede gelooss, fir mech virun d'Feier
aus Ranken ze setzen, dat am Häerd geknascht huet. Do hunn ech mech gutt lëschteg
gemaach iwwert de schaarfe \VUgras{Nordwand,} deen dat aalt \VUgras{Schëld}
\textsuperscript{(3)} knaupsen léisst an deen a schwaarze Wirbelen de
\VUgras{Rauch} \textsuperscript{(1)} ewechgedroen huet, dee aus dem Kamäin
erauskoum.

%%K t07-c1-04-1 p
4. --- Et war d'Stonn vum Mëttegiessen. Ouni méi laang ze waarden, hu mir dat
einfacht, mä schmackhaft Iesse gutt giess, dat een eis servéiert huet (*).

%%K t07-tit-3 sub
\VUpk{10.2pt}{40}{E puer Besuch.}
\VUsaut{3.0mm}

%%K t07-c1-05-1 p
5. --- Nom Iessen hunn ech mäi Papp begleet, deen Assureur ass, bei seng Clienten.
Fir d'éischt hu mir de \VUgras{Schmadd} \textsuperscript{(5)} besicht, deen nieft der
Auberge wunnt. Hie war amgaangen, e Päerd ze beschloen. Ech hunn d'Kraaft vu sengem
Gesell bewonnert, dee fest op senger Liederschäerz de \VUgras{Huff} vum Päerd
gehalen huet, während de Schmadd d'\VUgras{Eisen} \textsuperscript{(6)} mat staarke
Hummerschléi befestegt huet.

%%K t07-c1-06-1 p
6. --- Duerno si mir bei de \VUgras{Schouster} \textsuperscript{(7)} vum Duerf
gaangen, den alen Jhemp. Obwuel hien als Schëld eng wonnerschéin \VUgras{Bott} huet,
huet hie sech domat begnügt, en aalt Paar duerchgetruede Schong nei ze suelen.

%%K t07-c1-06-2 p
Den Alen huet eis laachend gesot: „An der Stad war ech „Schouster“ an hat keng
Clienten. Hei sinn ech „Schouflécker“, an et gëtt vill Aarbecht.“

%%K t07-c1-07-1 p
7. --- Hien huet eis vu sengem Noper, dem \VUgras{Coiffeur} \textsuperscript{(8)},
erzielt, deem seng Clienten net zefridde sinn. Seng \VUgras{Rasoiren} sinn ëmmer
schaarteg, a seng \VUgras{Schéiere} schneiden ni gutt.

Mä well hien deen eenzege Coiffeur vum Duerf ass, ass ee selbstverständlech
gezwongen, sech vun him raséieren an d'Hoer schneiden ze loossen. Iwwregens ass hien e
gäizegen an ongefällege Mann.

%%K t07-c1-08-1 p
8. --- Glécklecherweis sinn déi aner \VUgras{Noperen} him net ähnlech. Zum Beispill
ass de \VUgras{Stellmécher} \textsuperscript{(9)} e gudde Mann, fläisseg an
gefälleg. Et ass e Vergnügen, e Gefier vun him reparéieren ze loossen. Seng
Aarbecht ass ëmmer fäerdeg gemaach.

%%K t07-c1-09-1 p
9. --- Den alen Jhemp huet sech och net iwwert de \VUgras{Suedler}
\textsuperscript{(10)} bekloot, deem hien heiansdo hëlleft, d'\VUgras{Geschir} ze
neelen, wann d'Aarbecht dréngt.

%%K t07-c1-09-2 p
Mäi Papp huet e no senger Famill gefrot: „Alleguer si gesond. Meng Enkelin ass an der
\VUgras{Schoul} \textsuperscript{(12)}. Hir \VUgras{Léierin} \textsuperscript{(13)}
ass ganz zefridde mat hir, si ass eng gutt \VUgras{Schülerin} \textsuperscript{(14)}.
Si ass net wéi mäi Sprëtzert vu Jong; deen denkt nëmmen un d'Spillen. De
\VUgras{Léierer} \textsuperscript{(15)} bréngt him näischt bäi.“

%%K t07-tit-4 sub
\VUsaut{3.0mm}
\VUpk{10.2pt}{40}{Duerfgeschwätz.}
\VUsaut{3.0mm}

%%K t07-c1-10-1 p
10. --- Den Alen huet eis duerno d'Neiegkeete vum Duerf erzielt. Et gëtt eng nei
Schoul gebaut, well déi am \VUgras{Gemengenhaus} ass ze kleng ginn. Iwwregens ass dat
einfach, well et geet duer, en Deel vum Gaart vum \VUgras{Millermeeschter} ze kafen.
Dat wäert fir dee aarme Mann ganz virdeelhaft sinn. Wéinst der Keelt kënnen
d'\VUgras{Millesteng} vun der \VUgras{Millen} \textsuperscript{(16)} de Weess net méi
malen, well d'\VUgras{Rad} \textsuperscript{(17)} vum \VUgras{Äis}
\textsuperscript{(25)} vum \VUgras{Baach} \textsuperscript{(23)} agefruer
ugehale gouf.

%%K t07-c1-11-1 p
11. --- „Wat de Bau ubelaangt, hutt Dir eis nei \VUgras{Kierch}
\textsuperscript{(18)} gesinn? Si ass ganz molerësch ënner hirem Schnéimantel.
D'\VUgras{Kréien} \textsuperscript{(19)}, déi hiren alen Tuerm net méi erkennen,
sëtze traureg um grousse Bam vum \VUgras{Kierfecht} \textsuperscript{(20)}.“

%%K t07-c1-12-1 p
12. --- Dee Gedanken un de Kierfecht huet dem Schouster d'Leit erënnert, déi mäi
Papp kannt hat an déi d'lescht Joer gestuerwe sinn. „Wéi vill \VUgras{Griewer}
\textsuperscript{(21)}, mäi gudden Här, sinn zu deenen anere bäikomm ënnert der
\VUgras{Zypress} \textsuperscript{(22)}! Den Doud huet eisen alen Notaire ewechgeholl.
All d'Duerfleit waren op sengem \VUgras{Begriefnes} dobäi.“

%%K t07-c1-13-1 p
13. --- „Do ass säin Nofolger, dee um Baach Schlittschong leeft. Hien ass grad komm,
fir de Rimm vu sengem \VUgras{Schlittschong} \textsuperscript{(26)} bei mir
reparéieren ze loossen, deen ofgerappt war.“

%%K t07-c1-14-1 p
14. --- „Do ass och um Ufer vum Baach de neie \VUgras{Feldhidder}
\textsuperscript{(27)}. Dat ass en ale Gendarm, deen d'Bengele vum Duerf erféiert. Si
lafe gläich fort, wa si seng \VUgras{Guêtren} \textsuperscript{(29)} a säi
\VUgras{Kepi} \textsuperscript{(28)} gesinn.“

%%K t07-c1-15-1 p
15. --- „Ah! do ass den ale Josef, deen e \VUgras{Kärchen mat Holzbierdelen}
\textsuperscript{(30)} fir de Bäcker féiert. --- Dat ass wouer, sot mäi Papp, ech
erkennen säi Gespann vu \VUgras{Mailen} \textsuperscript{(31)}. Ech muss jo mat him
schwätzen, ech notzen d'Geleeënheet. Bis geschwënn, Papp Jhemp.“

%%K t07-c1-16-1 p
16. --- Grad wéi mir erausgaange sinn, hunn d'Kanner vum Duerf d'Schoul verlooss a
si lafend op d'\VUgras{Rutschbunn} \textsuperscript{(33)} gestierzt, mat der Gefor ze
falen a sech beim \VUgras{Fall} \textsuperscript{(34)} e Glidd ze briechen. Ee vun
hinnen huet säi \VUgras{Schlitt} \textsuperscript{(32)} beim Stellmécher geholl, huet
ee vu senge Kameraden dragesat, an ass lafend fortgaangen.

%%K t07-c1-17-1 p
17. --- Anerer si sech op e \VUgras{Schnéihaufen} \textsuperscript{(37)} nieft der
\VUgras{Bréck} \textsuperscript{(40)} gestierzt an hunn ugefaangen, déi
\VUgras{Schnéistatu} \textsuperscript{(39)} ze bombardéieren, déi si gëschter opgeriicht
haten an där si en Hutt, eng Päif an eng Schoultäsch iwwer der Schëller ugedoen haten.

%%K t07-c1-18-1 p
18. --- Si këmmere sech wéineg ëm de frostege Wand an d'Keelt. Déi meescht hunn net
emol eng \VUgras{Halsdicher} \textsuperscript{(35)}, a fir sech nom Kampf mat
\VUgras{Schnéiballen} \textsuperscript{(38)} nees ze wiermen, geet hinnen eng Handvoll
ganz waarm \VUgras{Kastanien} \textsuperscript{(42)} duer, déi si bei der aler
\VUgras{Verkeeferin} Jacobine kafen, déi ënner hirem ze dënne \VUgras{Schal}
\textsuperscript{(43)} vu Keelt zidderd.

%%K t07-tit-5 sub
\VUsaut{3.0mm}
\VUcentre{12.0pt}{0}{\textit{Zweet Zeen.}}
\VUsaut{3.0mm}

%%K t07-tit-6 sub
\VUpk{10.2pt}{40}{De Bauerenhaff am Fréijoer.}
\VUsaut{4.0mm}

%%K t07-c2-01-1 p
1. --- Trotz all de Freede vum Wanter begréisse mir mat déiwer Freed d'Erëmkomme vum
\VUgras{Fréijoer.}

%%K t07-c2-01-2 p
Wat e frëschent Bild ass e Bauerenhaff owes, am Mount Mee!

%%K t07-c2-02-1 p
2. --- Um Horizont leet d'\VUgras{Sonn} \textsuperscript{(1)} sech lues erof an
iwwerschwemmt d'\VUgras{Land} \textsuperscript{(3)} mat hire purpurroude
\VUgras{Strahlen} \textsuperscript{(2)}. De fläissege \VUgras{Bauer}
\textsuperscript{(6)} pressd sech, säin \VUgras{Acker} \textsuperscript{(4)} ze
plécken.

%%K t07-c2-02-2 p
Mat sengem \VUgras{Plou} \textsuperscript{(5)}, dee virun e kräftegt \VUgras{Päerd}
\textsuperscript{(7)} gespaant ass, zitt hien d'\VUgras{Fuerch}
\textsuperscript{(8)}, an déi de \VUgras{Séier} \textsuperscript{(9)} mat voller Hand
de \VUgras{Somen} geheit, deen déi gëllen Ähren ergëtt.

%%K t07-c2-03-1 p
3. --- D'\VUgras{Méier} \textsuperscript{(11)} mat hire Seesen hunn d'Gras vun der
Wiss gemäht, d'Hee-Wenderen hunn d'\VUgras{Hee} \textsuperscript{(10)} gedréchent,
andeems si et mat \VUgras{Forken} \textsuperscript{(12)} a Rechen ëmgedréint hunn, an
do kommen si zréck.

%%K t07-c2-03-2 p
Déi eng ginn virum schwéiere Won, deen vu Ochse gezunn gëtt; si droen hiert
Handwierksgeschir op der Schëller. Déi aner, méi bequem, hu sech gemittlech um
duftende Hee néiergelooss.

%%K t07-c2-04-1 p
4. --- Am Laanschtgoe gréisse si frëndlech de \VUgras{Gäertner}
\textsuperscript{(16)}, deen am \VUgras{Uebstgaart} \textsuperscript{(13)} amgaangen
ass, d'\VUgras{Uebstbeem} ze schneiden an d'\VUgras{Bierebeem}
\textsuperscript{(14)} ze pfropen. D'\VUgras{Quetschebeem} \textsuperscript{(15)} bléien
op, an am gutt gedéngte \VUgras{Geméisgaart} \textsuperscript{(17)} stinn d'Geméis,
d'\VUgras{Kabeskäpp} an d'\VUgras{Zalot} scho gutt.

%%K t07-c2-04-2 p
Op der klenger Mauer spreet e schéine \VUgras{Pauteil} \textsuperscript{(18)} säi
Schwanz aus, fir seng wonnerschéi Fiederen ze weisen.

%%K t07-tit-7 sub
\VUpk{10.2pt}{40}{De Bauerenhaff.}
\VUsaut{3.0mm}

%%K t07-c2-05-1 p
5. --- D'Dieren vum \VUgras{Päerdsstall} \textsuperscript{(19)} sinn op, an ee gesäit
d'Krëpp an d'\VUgras{Sträe} \textsuperscript{(23)} vun der \VUgras{Mär}
\textsuperscript{(21)}, déi de \VUgras{Päerdsknecht} \textsuperscript{(20)} grad
ausgespaant huet. Dat \VUgras{Fillen} \textsuperscript{(22)}, esou komesch mat senge
laange Been, leeft hipsend hannert senger Mamm hier.

%%K t07-c2-06-1 p
6. --- Nieft dem \VUgras{Pëtz} \textsuperscript{(29)} gi d'\VUgras{Kéi}
\textsuperscript{(25)} an den hëlzenen \VUgras{Trach} \textsuperscript{(26)} drénken,
deen hinnen als Drénkplaz déngt. Dat \VUgras{Kallef} \textsuperscript{(24)} notzt
d'Geleeënheet fir ze sëllen, an de \VUgras{Stéier} \textsuperscript{(27)}, deen de
gudde Geroch vum Fudder richt, brullt frou a hieft houfreg de Kapp mat de mächtege
\VUgras{Hierner} \textsuperscript{(28)}.

%%K t07-c2-07-1 p
7. --- All Leit schaffen. D'\VUgras{Déngschtmeedchen} \textsuperscript{(32)} hält an
der Hand de \VUgras{Seel} \textsuperscript{(31)} vum Pëtz. Si schäfft Waasser mat
engem \VUgras{Eemer} \textsuperscript{(30)}, fir d'Béischten ze drénken ze ginn. Nieft
der \VUgras{Scheier} \textsuperscript{(33)} rechent e Mann d'Stréi, a virun der Dier
vum \VUgras{Haus} \textsuperscript{(34)} spënnt eng al \VUgras{Spënnerin}
\textsuperscript{(35)} mat hirem Spannrad an hirem \VUgras{Rocken} de
\VUgras{Läin} oder den \VUgras{Hanf} wéi an aler Zäit.

%%K t07-tit-8 sub
\VUsaut{3.0mm}
\VUpk{10.2pt}{40}{De Bauer.}
\VUsaut{3.0mm}

%%K t07-c2-08-1 p
8. --- De \VUgras{Bauer} \textsuperscript{(36)} leet all dës Aarbechten a gëtt senge
Kniechten Uweisungen. Hie recommandéiert, d'\VUgras{Eg} \textsuperscript{(40)} an
d'\VUgras{Haw} \textsuperscript{(39)} ënner Daach ze setzen, déi ee bei der
\VUgras{Hütt} \textsuperscript{(38)} vergiess huet, an där den \VUgras{Hond}
\textsuperscript{(37)} wunnt.

%%K t07-c2-09-1 p
9. --- Säi Rufe schreckt eng \VUgras{Dauf} \textsuperscript{(41)} op, déi mat alle
Flilleke bis an den \VUgras{Dauwenhaus} \textsuperscript{(42)} flitt, an
d'\VUgras{Hénger} \textsuperscript{(43)}, déi sech pressen, an den
\VUgras{Hénkelhaff} \textsuperscript{(44)} zréckzegoen.

%%K t07-c2-10-1 p
10. --- Virum \VUgras{Schofstall} \textsuperscript{(48)}, do ass den ale
\VUgras{Schäfer} \textsuperscript{(45)}, deen seng \VUgras{Herd}
\textsuperscript{(49)} zréckféiert. Andeems hie säi \VUgras{Stack}
\textsuperscript{(46)} hält, dee him ni feelt, sou wéi seng ausgeblatschten
\VUgras{Bluse} \textsuperscript{(47)}, féiert hien an de Stall d'\VUgras{Hämmelen}
\textsuperscript{(50)}, d'\VUgras{Schof} \textsuperscript{(53)}, d'\VUgras{Lämmer}
\textsuperscript{(54)}, d'\VUgras{Geessen} \textsuperscript{(51)} an d'\VUgras{Zicklein}
\textsuperscript{(52)}. Säin trei \VUgras{Hond} \textsuperscript{(55)} Argus kënnt als
Leschten an dreift mat sengem Gebell déi Nozügler un.

%%K t07-c2-11-1 p
11. --- All dee Kaméidi an dat Gerappels stéiert net d'Rou vun der gudder
\VUgras{Mëllechkou} \textsuperscript{(56)}, déi hir \VUgras{Klack}
\textsuperscript{(57)} klengele léisst, während ee si virun der Dier vum
\VUgras{Stall} \textsuperscript{(58)} mëllecht.

%%K t07-c2-12-1 p
12. Si schéngt ze verstoen, datt si hir Mëllech muss ginn, fir datt d'\VUgras{Bäuerin}
\textsuperscript{(60)} \VUgras{Botter} an der \VUgras{Botterkiirn}
\textsuperscript{(61)} maache kann oder déi ganz gutt \VUgras{Kéis}
\textsuperscript{(64)}, déi vun deem Brett drëpsen, dat um \VUgras{Bidden}
\textsuperscript{(63)} läit.

%%K t07-c2-13-1 p
13. --- Déi ganz \VUgras{Mëllecherei} \textsuperscript{(59)} gëtt ganz propper vum
\VUgras{Bauerenknecht} \textsuperscript{(65)} gehalen, deen d'\VUgras{Mëllechdëppen}
\textsuperscript{(62)} grad am grousse Eemer gewäsch huet, deen hien an der Hand
dréit.

%%K t07-tit-9 sub
\VUsaut{3.0mm}
\VUpk{10.2pt}{40}{Den Hénkelhaff.}
\VUsaut{3.0mm}

%%K t07-c2-14-1 p
14. --- Mat senge décken Holzschong geet hien e bësse schwéierfälleg, an esou jagt
hien de \VUgras{Pouler} \textsuperscript{(68)}, d'\VUgras{Enten}
\textsuperscript{(66)}, d'\VUgras{Enterecher} \textsuperscript{(67)} an
d'\VUgras{Gäns} \textsuperscript{(69)} op, déi hire \VUgras{Schniewel}
\textsuperscript{(70)} wäit opmaachen an onrouheg jäizen.

%%K t07-c2-14-2 p
Den \VUgras{Hunn} \textsuperscript{(71)}, méi kämpferesch, riicht säi roude
\VUgras{Kamm} \textsuperscript{(72)} op, rëselt d'Fiedere vu sengem
\VUgras{Schwanz} \textsuperscript{(73)} a léisst e lauden „kikeriki“ héieren.

%%K t07-c2-15-1 p
15. --- D'\VUgras{Gluckhong} \textsuperscript{(74)} sicht Fudder um \VUgras{Mëscht}
\textsuperscript{(81)} mat senge \VUgras{Kichelcher} \textsuperscript{(75)}. Dat
\VUgras{Ferkel} \textsuperscript{(78)} leeft bei seng Mamm, déi enorm \VUgras{Sau}
\textsuperscript{(77)}, déi mir nieft dem Saustall gesinn.

%%K t07-c2-16-1 p
16. --- An hire Käfeger iessen d'\VUgras{Kanéngercher} \textsuperscript{(79)}
gierdeg dat zaart \VUgras{Gras} \textsuperscript{(80)}, dat hinnen d'Duechter vum
Bauer bréngt. D'Jeanne huet hir Kanéngercher ganz gär, an all Dag, wa si déi frësch
\VUgras{Eeër} \textsuperscript{(76)} am Hénkelhaff geholl huet, kënnt si hir
Frëndercher besichen.

\VUsaut{6.0mm}
\VUfilet{22mm}
